% ============================================================================
% SEC03.TEX - Section 3.3: Variabel dan Tipe Data
% ============================================================================

\section{Variabel dan Tipe Data}

\begin{learningoutcomes}
\item Menggunakan tipe data dasar C\#
\item Memahami tipe data khusus Unity (Vector3, Quaternion, dll)
\item Menggunakan public dan private variables
\item Memahami SerializeField
\end{learningoutcomes}

\subsection{Tipe Data Dasar}

\begin{itemize}
\item \textbf{int}: Bilangan bulat
\item \textbf{float}: Bilangan desimal
\item \textbf{bool}: True/false
\item \textbf{string}: Teks
\item \textbf{char}: Satu karakter
\end{itemize}

\subsection{Tipe Data Unity}

\begin{itemize}
\item \textbf{Vector3}: Vektor 3D (x, y, z)
\item \textbf{Vector2}: Vektor 2D (x, y)
\item \textbf{Quaternion}: Rotasi 3D
\item \textbf{Color}: Warna RGBA
\item \textbf{GameObject}: Reference ke GameObject
\item \textbf{Transform}: Reference ke Transform component
\end{itemize}

\subsection{Contoh Implementasi}

\begin{lstlisting}[language={[Sharp]C}]
using UnityEngine;

public class VariableExample : MonoBehaviour
{
    // Public: Terlihat di Inspector
    public int health = 100;
    public float speed = 5.5f;
    public string playerName = "Player1";
    public Vector3 position;
    
    // Private: Tidak terlihat di Inspector
    private int score;
    
    // SerializeField: Private tapi terlihat di Inspector
    [SerializeField] private int maxHealth = 100;
}
\end{lstlisting}

\begin{catatan}
Public variables muncul di Inspector dan dapat diubah tanpa mengedit kode. Gunakan SerializeField untuk private variables yang perlu diatur di Inspector.
\end{catatan}
